\section{The Real Numbers}
\label{sec:real-numbers}

\subsection{Defining the Real Numbers}
\label{sec:defining-real-numbers}

Recall that we defined \(\NN\) using the Peano's Axioms, and we may obtain \(\ZZ\) from \(\NN\) by allowing for 'subtraction'.

\begin{definition}[\(\ZZ\)]
    Formally, \(\ZZ\) is the set of equivaleence classes of \(\NN \times \NN\) under the equivalence relation
    \[
        \para{a, b} \sim \para{c, d} \iff a + d = b + c.
    \]
\end{definition}

We may think of \(\para{a, b}\) as \(a - b\). We denote
\[
    0 = \brac{\para{1, 1}}
\]
and
\[
    -a = \brac{\para{1, a + 1}}.
\]

We may similarly define addition and multiplication on \(\times\) using this pair, as
\[
    \para{a, b} + \para{c, d} = \para{a + c, b + d}
\]
and
\[
    \para{a, b} \times \para{c, d} = \para{ac + bd, bc + ad},
\]
and we may indeed check that these operations are well-defined.

Similarly, we may obtain \(\QQ\) from \(\ZZ\) by allowing for 'division'.

\begin{definition}[\(\QQ\)]
    Formally, \(\QQ\) is the set of the equivalence classes of \(\ZZ \times \NN\) under the equivalence relation
    \[
        \para{a, b} \sim \para{c, d} \iff ad = bc.
    \]
\end{definition}

Similarly, notational wise, we write \(\frac{a}{b}\) instead of \(\brac{\para{a, b}}\), and we may similarly define \(+\) and \(\times\) by
\[
    \para{a, b} + \para{c, d} = \para{ad + bc, bd}
\]
and
\[
    \para{a, b} \times \para{c, d} = \para{a \times c, b \times d}.
\]

In addition to this, we may define an \emph{ordering} on \(\QQ\), \(<\), which satisfy the following properties, that
\begin{enumerate}
    \item if \(a, b \in \QQ\), then precisely one of \(a < b, a = b, a > b\) holds;
    \item if \(a < b\) and \(b < c\), then \(a < c\).
\end{enumerate}

This ordering is useful, in the sense that there is a rational number between any two rational numbers. If \(p, q \in \QQ\) with \(p < q\), w e have
\[
    p < \frac{p + q}{2} < q
\]
so \(\QQ\) is \emph{densely ordered}.

However, this does not ensure that all the gaps are filled in.

\begin{proposition}[\(\sqrt{2}\) is irrational]
    There is no rational \(x\) with \(x^2 = 2\).    
\end{proposition}

\begin{proof}[by Number Theory]
    Suppose that \(x^2 = 2\). Note we may assume \(x > 0\) since \(\para{-x}^2 = x^2\). If \(x\) is rational, then
    \[
        x = \frac{a}{b}
    \]
    for some \(a, b \in \NN\). Thus, \(a^2 = 2b^2\).

    But note that the exponent of \(2\) in the prime factorisation on the left-hand side is even, while the exponent of \(2\) in the prime factorisation on the right-hand side is odd. This contradicts with the fundamental theorem of arithemetics.
\end{proof}

\begin{remark}
    This proof reflects from a number theory perspective, that if there is \(x \in \QQ\) with \(x^2 = n\) for some \(n \in \NN\), then \(n\) must itself be a square number.
\end{remark}

\begin{proof}[by Algebra]
    Suppose \(x^2 = 2\) for some \(x = \frac{a}{b}\), with \(a, b \in \NN\). Then, for any \(c, d \in \ZZ\), \(\para{cx + d}\) is of the form \(\frac{e}{b}\) for some \(e \in \ZZ\).

    Then, if \(cx + d > 0\), we have \(cx + d \geq \frac{1}{b}\).

    But \(0 < x - 1 < 1\) as \(1 < x < 2\), if \(n\) is sufficiently large, we have
    \[
        0 < \para{x - 1}^n < \frac{1}{b}.
    \]

    But for any \(n \in \NN\), \(\para{x - 1}^n\) is of the form \(cx + d\), for some \(c, d \in \ZZ\), since \(x^2 = 2\). This leads to a contradiction.
\end{proof}

\begin{remark}
    This proof shows that if the square root of an integer is rational, then it must in fact be a whole number.
\end{remark}

This shows that \(\QQ\) has a gap. But how do we fill in this gap only making reference to \(\QQ\), without mentioning \(\RR\)? (We haven't defined the real numbers yet.)

\(\QQ\) fails another property which is key to constructing the reals. Let \(A\) be the set of positive rationals \(p\), such that \(p^2 < 2\). Formally, let
\[
    A = \set{p \in \QQ}{p^2 < 2}.
\]

We show that \(A\) has no largest number. For any \(p \in A\), consider \(q = p - \frac{p^2 - 2}{p + 2}\).

Since \(p \in A\), \(p^2 - 2 < 0\), so we have \(q > p\). Also,
\[
    q^2 - 2 = \frac{2 \para{p^2 - 2}}{\para{p + 2}^2} < 0,
\]
so \(q \in A\).

Similarly, the set
\[
    \set{q \in \QQ}{q > 0, q^2 > 0}
\]
has no smallest element.

The \textbf{crucial point} is not about the \emph{maximum/minimum} of a set -- it is about a \emph{least upper bound} -- which is the "gap". (And luckily, all the "gaps".)

\begin{definition}[\(\RR\)]
    The \emph{real numbers} \(\RR\) are a set with elements \(0, 1 \in \RR\) where \(0 \neq 1\), equipped with operations \(+\) and \(\times\), and an ordering \(<\) satisfying the following axioms:
    \begin{enumerate}
        \item \(\para{\RR, +, 0}\) is an abelian group: \(+\) is commutative and associative with identity \(0\), and inverses exist under \(+\);
        \item \(\para{\RR\multiplicative, \times, 1}\) is an abelian group: \(\times\) is commutative and associative with identity \(1\), and inverses exist under \(\times\) for non-zero elements;
        \item \(\times\) is distributive over \(+\):
        \[
            a \times \para{b + c} = a \times b + a \times c,
        \]
        and \(\para{\RR, +, \times}\) form a field;
        \item \(\para{\RR, <}\) is a totally ordered field: for all \(a, b\), exactly one of \(a < b\), \(a = b\) and \(a > b\) holds; for all \(a, b, c\), if \(a < b\) and \(b < c\), then \(a < c\);
        \item the ordering preserve arithemetics: for all \(a, b, c\), if \(a < b\), then \(a + c < b + c\), and \(ac < bc\) for \(c > 0\);
        \item given any set \(S\) of reals that is non-empty and bounded above, \(S\) has a \emph{least upper bound}.
    \end{enumerate}

    The final axiom is known as the \emph{least upper bound axiom.}
\end{definition}

\begin{definition}[Least Upper Bound]
    We say a set \(S\) is \emph{bounded above} if there is some \(x \in \RR\), such that \(x \geq y\) for all \(y \in S\). Such an element \(x\) is called an \emph{upper bound of \(S\)}.

    We say \(x\) is \emph{the least upper bound} for \(S\), if \(x\) is an upper bound for \(S\), and any other upper bound \(x'\) satisfies \(x' \geq x\).

    We denote this as \(x = \LUB\para{S}\), \(x = \supremum\para{S}\), or most commonly,
    \[
        x = \sup\para{S}.
    \]
\end{definition}

Note that we defined \(\RR\) -- but its existence is not immediately obvious from the definition. This can be seen from a technique called \emph{Dedkind Cuts}.

\begin{remarks}
    \begin{enumerate}
        \item From the axioms, we can check, for example, that \(0 < 1\). Indeed, if not, then \(1 < 0\). Since \(0 \neq 1\), assume on the countrary that \(1 < 0\), so we have
        \[
            0 = 1 - 1 < 0 - 1 = -1,
        \]
        so
        \[
            0 = 0 \cdot \para{-1} < \para{-1} \para{-1} = 1
        \]
        which gives a contradiction.
        \item We may consider \(\QQ\) as \emph{contained} in \(\RR\), by identifying \(\frac{a}{b} \in \QQ\) with \(a \cdot b\inverse \in \RR\).
        \item \(\QQ\) satisfies all axioms \emph{apart from} the final one -- the least upper bound axiom. For example, the set \(A\) involving squaring to \(2\) above does not have a supremum.
        \item The words non-empty and bounded above are crucial.
        
        If \(S\) is empty, then every \(x \in \RR\) is trivially an upper bound, so there is no least upper bound.

        If \(S\) is not bounded above, then it has no upper bound, and certainly does not have a least one.
    \end{enumerate}
\end{remarks}

We give some examples on the idea of the least upper bound.

\begin{examples}
    \begin{enumerate}
        \item Consider \(S = \set{x \in \RR}{0 \leq x \leq 1} = \ccintv{0}{1}\).
        
        \(2\) is an upper bound for \(S\), since for all \(x \in S\), \(x \leq 2\).

        \(\frac{3}{4}\) is \emph{not} an upper bound for \(S\), since \(\frac{7}{8} \in S\), and \(\frac{7}{8} > \frac{3}{4}\).

        In fact, we have
        \[
            \sup \para{S} = 1,
        \]
        since \(1\) is an upper bound: for all \(x \in S, x \leq 1\); and every other upper bound is \(\geq 1\), since \(1 \in S\).

        \item Consier \(S = \set{x \in \RR}{0 < x < 1} = \oointv{0}{1}\).
        
        Similarly, \(2\) is an upper bound for \(S\), and \(\frac{3}{4}\) is not.

        In fact, we have
        \[
            \sup \para{S} = 1,
        \]
        since \(1\) is an upper bound: for all \(x \in S, x < 1\) so \(x \leq 1\); and no opper bound \(c\) satisfies \(c < 1\).

        Indeed, \(c\) is certainly greater than \(0\), and furthermore \(c \geq \frac{1}{2}\) since \(\frac{1}{2} \in S\).

        So if \(c < 1\), then \(0 < c < 1\), so \(\frac{c + 1}{2} \in S\), but we have \(\frac{c + 1}{2} > c\).
    \end{enumerate}
\end{examples}

\begin{remark}
    If \(S\) have a greatest element, then
    \[
        \sup \para{S} = \max \para{S} \in S.
    \]

    However, as the example shows, \(\sup\para{S}\) can exist when \(\max\para{S}\) does not: in this case, \(\sup\para{S} \notin S\).
\end{remark}

Now we want to find the least upper bound of
\[
    S = \brce{0, \frac{1}{2}, \frac{2}{3}, \frac{3}{4}, \ldots} = \set{1 - \frac{1}{n}}{n \in \NN}.
\]

Clearly, \(1\) is an upper bound of \(S\). But is there an upper bound less than 1?

\begin{proposition}[Axiom of Archimedes]
    \(\NN\) is not bounded above in \(\RR\).
\end{proposition}

\begin{proof}
    Suppose on the contrary, that \(\NN\) is bounded above, then it has a least upper bound
    \[
        c = \sup \para{\NN} \in \RR.
    \]

    By definition, \(\para{c - 1}\) is not an upper bound for \(\NN\), so there is some \(n \in \NN\), such that \(n > c - 1\).

    Then, \(n + 1 \in \NN\), with \(n + 1 > c\), contradicting that \(c\) is an upper bound for \(\NN\).
\end{proof}

\begin{corollary}[Reciprocal of Natural Numbers can be Arbitrarily Small]
    For any real number \(t > 0\), there exists \(n \in \NN\) such that \(\frac{1}{n} < t\).
\end{corollary}

\begin{proof}
    Given \(t > 0\), by Axiom of Archimedes, there is \(n \in \NN\) where \(n > \frac{1}{t}\).

    Hence, \(\frac{1}{n} < t\) as desired.
\end{proof}

\begin{definition}[Greatest Lower Bound]
    A set \(S\) is said to be \emph{bounded below}, if there exists \(x\), such that for all \(y \in S\), we have \(x \leq y\). Such an \(x\) is called a \emph{lower bound} for \(S\).

    We say \(x\) is the \emph{greatest lower bound} for \(S\), if \(x\) is a lower bound for \(S\), and any other lower bound \(x'\) satisfies \(x' \leq x\).

    We denote this as \(x = \GLB \para{S}\), \(x = \infimum\para{S}\), or most commonly,
    \[
        x = \inf \para{S}.
    \]
\end{definition}

Note that if \(S\) is empty and bounded below, then
\[
    -S = \set{-y}{y \in S}
\]
is non-empty and bounded above, so it has a least upper bound
\[
    \sup \para{-S} = c,
\]
and hence \(-c\) is the greatest lower bound.

The previous corollary immediately implies that
\[
    \inf \para{\set{\frac{1}{n}}{n \in \NN}} = 0.
\]

The previous proposition and corollary can be interpreted as that there are no 'infinitely large' or 'infinitely small' numbers in \(\RR\).

Now, we return to a previous example.

\begin{example}
    We want to find the least upper bound of
    \[
        S = \brce{0, \frac{1}{2}, \frac{2}{3}, \frac{3}{4}, \ldots} = \set{1 - \frac{1}{n}}{n \in \NN}.
    \]

    We claim \(\sup \para{S} = 1\), and suppose on the countrary that \(c < 1\) is an upper bound for \(S\). Then \(1 - \frac{1}{n} \leq c\), for all \(n \in \NN\), so
    \[
        0 < 1 - c \leq \frac{1}{n}
    \]
    for all \(n \in \NN\), contradicting with the previous corollary.
\end{example}

Now we show that the gaps we found in \(\QQ\) beforehand is filled by the least upper bound axiom.

\begin{theorem}[Square Roots Exist in \(\RR\)]
    There exists \(x \in \RR\) with \(x^2 = 2\).
\end{theorem}

\begin{proof}
    Let
    \[
        S = \set{x \in \RR}{x^2 < 2}.
    \]

    First we note that \(S\) is non-empty since \(1 \in S\), and it is bounded above by \(2\)

    Hence, \(S\) has a supremum, and let \(c = \sup S\). We observe \(1 < c < 2\), ad we claim that \(c^2 = 2\).

    Suppose that \(c^2 < 2\). For \(0 < t < 1\), we have
    \begin{align*}
        \para{c + t}^2 &= c^2 + 2ct + t^2 \\
        &< c^2 + 4t + t \\
        &= c^2 + 5t \\
        &< 2,
    \end{align*}
    for sufficiently small \(t\), say \(t < \frac{2 - c^2}{5}\).

    But this contradicts with the assumption that \(c\) is \emph{an upper bound} for \(S\), since \(c + t \in S\) and \(c + t > c\).

    Now, suppose that \(c^2 > 2\). For \(0 < t < 1\), we have
    \begin{align*}
        \para{c - t}^2 &= c^2 - 2ct + t^2 \\
        &> c^2 - 4t + 0 \\
        &> 2,
    \end{align*}
    for sufficiently small \(t\), say \(t < \frac{c^2 - 2}{4}\).

    This contradicts with the assumption that \(c\) is the \emph{least} upper bound of \(S\), since \(c - t < c\) and \(c - t\) is an upper bound for \(S\).

    Hence, \(c^2 = 2\).
\end{proof}

\begin{remark}
    A similar proof shows that \(\sqrt[n]{x}\) exists for all \(n \in \NN\), for all \(x \in \RR\) where \(x > 0\):
    \[
        \forall n \in \NN, \forall x \in \oointv{0}{\plusinfty}, \exists y \in \RR: y^n = x.
    \]
\end{remark}

\begin{definition}[Irrational Numbers]
    A real that is not rational is \emph{irrational}.
\end{definition}

\begin{examples}[Examples of Irrationals]
    \begin{enumerate}
        \item \(\sqrt{2}, \sqrt{3}, \sqrt{5}, \sqrt{6}, \ldots\) are irrational.
        \item If \(2 + 3 \sqrt{5} = \frac{a}{b}\) for some \(a, b \in \NN\), then
        \[
            \sqrt{5} = \frac{a - 2b}{3b} \in \QQ
        \]
        but \(\sqrt{5}\) is irrational. So \(2 + 3 \sqrt{5}\) is irrational.
    \end{enumerate}
\end{examples}

\begin{proposition}[\(\QQ\) is dense in \(\RR\)]
    The rationals are \emph{dense} in the reals. Formally,
    \[
        \forall a, b \in \RR, a < b, \exists c \in \QQ: a < c < b.
    \]
\end{proposition}

\begin{proof}
    We may assume \(a \geq 0\). By corollary of the axiom of achimedes, we have some \(n \in \NN\) such that \(\frac{1}{n} < b - a\).

    Let
    \[
        T = \set{k \in \NN}{\frac{k}{n} \geq b}.
    \]

    By axiom of archimedes, there is some \(N \in \NN\) such that \(N > b\).
    
    Hence, we have
    \[
        \frac{Nn}{n} \geq b,
    \]
    and hence \(Nn \in T\), so \(T \neq \emptyset\).

    By the well-ordering principle, \(T\) has some least element \(m\). Now, let \(c = \frac{m - 1}{n} \in \QQ\).

    Since \(m - 1 \notin T\), we have \(c < b\).

    Now, if \(c \leq a\), then
    \[
        \frac{m}{n} = c + \frac{1}{n} < a + \para{b - a} = b,
    \]
    contradicts with \(m \in T\).

    Therefore, \(c > a\), and hence \(a < c < b\) with \(c \in \QQ\), as desired.
\end{proof}

\begin{proposition}[\(\RR \setminus \QQ\) is dense in \(\RR\)]
    The irrationals are \emph{dense} in the reals. Formally,
    \[
        \forall a, b \in \RR, a < b, \exists c \in \RR \setminus \QQ: a < c < b.
    \]
\end{proposition}

\begin{proof}
    Indeed, take \(c \in \QQ\) such that
    \[
        a \sqrt{2} < c < b \sqrt{2}
    \]
    then
    \[
        a < \frac{c}{\sqrt{2}} < b.
    \]
\end{proof}

\subsection{Sequences}
\label{sec:sequences}